\documentclass[10pt,a4paper]{article}
\input{../00_COMPILACION/t04_estilo_v2.tex}
\begin{document}
\tituloT{Tema 4. Números enteros. Divisibilidad. Números primos. Congruencia}{T04.02 · Tema de examen modular · versión 4 canónica · núcleo A/B y ampliaciones C}

\begin{idea}[Regla de escritura]
El texto ordinario constituye la columna vertebral. Los recuadros C pueden omitirse completos sin romper el hilo. La versión se cierra siempre con proyección educativa, conclusión y bibliografía; nunca se sacrifica el núcleo para conservar una ampliación.
\end{idea}

\section*{Índice estratégico}
\begin{tabularx}{\textwidth}{C{0.75cm}Y}
\toprule
1 & Introducción.\\
2 & Enteros y división euclídea. Construcción formal de $\Z$ (C).\\
3 & Divisibilidad, MCD, MCM, Bézout y ecuaciones diofánticas.\\
4 & Primos y factorización. Infinitud y funciones aritméticas (C).\\
5 & Congruencias, unidades, cancelación y congruencias lineales.\\
6 & Euler--Fermat y teorema chino del resto. Wilson (C).\\
7 & Aplicaciones y proyección educativa. RSA (C).\\
8 & Conclusión, bibliografía y normativa.\\
\bottomrule
\end{tabularx}

\section{Introducción}
Los números enteros constituyen el primer sistema numérico en el que la resta es una operación interna. En ellos surgen dos preguntas estructurales: cuándo un entero divide a otro y qué información permanece al considerar solamente los restos. De ambas nacen la divisibilidad, la teoría de números primos y las congruencias, que transforman problemas potencialmente ilimitados en problemas sobre factores o clases residuales.

En los libros VII y IX de los \emph{Elementos}, Euclides sistematizó el algoritmo del máximo común divisor y la teoría elemental de los primos. Boyer y Merzbach muestran que estas referencias responden a problemas matemáticos concretos: comparar medidas y descomponer números. Mucho después, Gauss convirtió las congruencias en un cálculo aritmético autónomo.

Partiremos de la estructura de $\Z$ y de la división euclídea; enlazaremos el algoritmo de Euclides con Bézout y las ecuaciones diofánticas; estudiaremos la factorización prima; y culminaremos con el cálculo modular, Euler--Fermat, el teorema chino del resto y sus aplicaciones.

\section{Enteros y división euclídea}
\begin{ampliacion}[· Construcción formal de $\Z$]
En $\N^2$ se define $(a,b)\sim(c,d)$ si $a+d=b+c$. La clase $[(a,b)]$ representa la diferencia $a-b$, y el cociente $\N^2/\!\sim$ se identifica con $\Z$. Se definen
\[
[(a,b)]+[(c,d)]=[(a+c,b+d)],\qquad
[(a,b)]\,[(c,d)]=[(ac+bd,ad+bc)].
\]
Las operaciones están bien definidas respecto de la equivalencia. Así, $(\Z,+,\cdot)$ es un anillo conmutativo con unidad y dominio de integridad; sus unidades son $\pm1$. El orden es compatible con la suma y con el producto por enteros positivos.
\end{ampliacion}

En adelante basta recordar que $(\Z,+,\cdot)$ es un dominio de integridad totalmente ordenado y que sus unidades son $\pm1$.

\begin{teorema}[División euclídea]
Dados $a\in\Z$ y $b\ne0$, existen únicos $q,r\in\Z$ tales que
\[
a=bq+r,\qquad 0\le r<|b|.
\]
\end{teorema}
\begin{proof}
Para $b>0$, el conjunto $S=\{a-bk:k\in\Z,\ a-bk\ge0\}$ es no vacío. Sea $r=a-bq$ su mínimo. Si $r\ge b$, entonces $r-b=a-b(q+1)$ sería un elemento menor de $S$; luego $0\le r<b$. Si hubiera dos expresiones, $b(q-q')=r'-r$ y $|r'-r|<b$, de donde $r=r'$ y $q=q'$. Para $b<0$ se emplea $|b|$.
\end{proof}
Por ejemplo, $-47=6(-8)+1$; exigir $0\le r<|b|$ impide tomar el resto $-5$.

\section{Divisibilidad, MCD, MCM y ecuaciones diofánticas}
Para $a,b\in\Z$, se escribe $a\mid b$ cuando existe $k\in\Z$ con $b=ak$. Si $a\mid b$ y $a\mid c$, entonces $a\mid(ub+vc)$ para cualesquiera $u,v\in\Z$.

Para $a,b$ no ambos nulos, el máximo común divisor positivo $d=\mcd(a,b)$ divide a ambos y todo divisor común divide a $d$. Si $a=bq+r$, los pares $(a,b)$ y $(b,r)$ poseen los mismos divisores comunes, por lo que
\[
\mcd(a,b)=\mcd(b,r).
\]
La iteración constituye el algoritmo de Euclides y termina porque los restos no negativos decrecen estrictamente.

\begin{teorema}[Identidad de Bézout]
Si $d=\mcd(a,b)$, existen $u,v\in\Z$ tales que $d=au+bv$.
\end{teorema}
\begin{proof}
El último resto no nulo del algoritmo es $d$. Sustituyendo hacia atrás cada resto en función de los anteriores se obtiene $d=au+bv$. Recíprocamente, todo divisor común de $a$ y $b$ divide cualquier combinación $au+bv$.
\end{proof}
Para $252$ y $198$ se obtiene $18=4\cdot252-5\cdot198$.

Si $a=\prod p_i^{\alpha_i}$ y $b=\prod p_i^{\beta_i}$, el MCD toma los exponentes mínimos y el MCM los máximos. Como
\[
\min(\alpha_i,\beta_i)+\max(\alpha_i,\beta_i)=\alpha_i+\beta_i,
\]
se deduce $\mcd(a,b)\mcm(a,b)=|ab|$ cuando $ab\ne0$.

La ecuación diofántica $ax+by=c$ tiene solución entera si y solo si $d\mid c$. Si $(x_0,y_0)$ es una solución particular, todas son
\[
x=x_0+\frac{b}{d}t,\qquad y=y_0-\frac{a}{d}t,\qquad t\in\Z.
\]
Hunacek presenta división euclídea, Euclides, Bézout y ecuaciones diofánticas como una única cadena conceptual, no como reglas aisladas.

\section{Números primos y factorización única}
Un entero $p>1$ es primo si sus únicos divisores positivos son $1$ y $p$; $0$ y $1$ no son primos.

\begin{lema}[Euclides]
Si $p$ es primo y $p\mid ab$, entonces $p\mid a$ o $p\mid b$.
\end{lema}
\begin{proof}
Si $p\nmid a$, entonces $\mcd(p,a)=1$. Por Bézout existen $u,v$ con $up+va=1$. Multiplicando por $b$, los dos sumandos del miembro izquierdo son divisibles por $p$, luego $p\mid b$.
\end{proof}

\begin{teorema}[Teorema fundamental de la aritmética]
Todo entero $n>1$ es producto de primos de forma única, salvo el orden.
\end{teorema}
La existencia se prueba por inducción fuerte. Para la unicidad, el lema de Euclides permite identificar y cancelar sucesivamente los factores primos.

\begin{ampliacion}[· Infinitud de los primos]
Si $p_1,\ldots,p_k$ fueran todos los primos, el número $N=p_1\cdots p_k+1$ no sería divisible por ninguno de ellos; sin embargo, $N>1$ posee un divisor primo, contradicción.
\end{ampliacion}

\begin{ampliacion}[· Funciones aritméticas]
Si $n=\prod p_i^{\alpha_i}$, entonces
\[
\tau(n)=\prod(\alpha_i+1),\qquad
\sigma(n)=\prod\frac{p_i^{\alpha_i+1}-1}{p_i-1}.
\]
Por ejemplo, $360=2^3\cdot3^2\cdot5$ y $\tau(360)=24$. Apostol desarrolla $\tau$, $\sigma$ y $\varphi$ como funciones multiplicativas.
\end{ampliacion}

\section{Congruencias, unidades y ecuaciones lineales}
Para $m\ge2$ se define
\[
a\equiv b\pmod m\quad\Longleftrightarrow\quad m\mid(a-b).
\]
Es una relación de equivalencia compatible con suma y producto. Sus clases forman el anillo $\Z/m\Z$. Una clase $\overline a$ es invertible exactamente cuando $\mcd(a,m)=1$.

\begin{proposicion}[Caso cuerpo]
El anillo $\Z/m\Z$ es cuerpo si y solo si $m$ es primo.
\end{proposicion}
\begin{proof}
Si $m$ es primo y $\overline a\ne\overline0$, Bézout proporciona un inverso de $\overline a$. Si $m=rs$ es compuesto, con $1<r,s<m$, entonces $\overline r$ y $\overline s$ son no nulas pero su producto es cero.
\end{proof}

La cancelación exige controlar las hipótesis. En general,
\[
ac\equiv bc\pmod m\quad\Longrightarrow\quad a\equiv b\pmod{m/\mcd(c,m)}.
\]
Por ejemplo, $2\cdot1\equiv2\cdot4\pmod6$, aunque $1\not\equiv4\pmod6$.

\begin{teorema}[Congruencia lineal]
La ecuación $ax\equiv b\pmod m$ tiene solución si y solo si $d=\mcd(a,m)$ divide a $b$; en ese caso posee exactamente $d$ soluciones incongruentes módulo $m$.
\end{teorema}
Equivale a la ecuación diofántica $ax-my=b$. Por ejemplo, $14x\equiv8\pmod{20}$ se reduce a $7x\equiv4\pmod{10}$ y produce $x\equiv2,12\pmod{20}$. Ivorra conecta de forma natural Bézout, inversos y congruencias lineales.

\section{Euler--Fermat y teorema chino del resto}
Sea $\varphi(m)$ el número de clases invertibles módulo $m$.

\begin{teorema}[Euler--Fermat]
Si $\mcd(a,m)=1$, entonces $a^{\varphi(m)}\equiv1\pmod m$. En particular, si $p$ es primo, $a^p\equiv a\pmod p$.
\end{teorema}
Multiplicar por $a$ permuta las clases invertibles; al comparar sus productos y cancelar se obtiene el resultado.

\begin{teorema}[Teorema chino del resto]
Si $m_1,\ldots,m_r$ son coprimos dos a dos, el sistema $x\equiv a_i\pmod{m_i}$ tiene una única solución módulo $M=m_1\cdots m_r$.
\end{teorema}
Con $M_i=M/m_i$ y $u_iM_i\equiv1\pmod{m_i}$, una solución es
\[
x_0=\sum_{i=1}^r a_iM_iu_i.
\]
La unicidad se sigue porque dos soluciones difieren en un múltiplo de cada $m_i$ y, por coprimalidad, de $M$.

\begin{ampliacion}[· Wilson]
Un entero $p>1$ es primo si y solo si $(p-1)!\equiv-1\pmod p$. Aporta amplitud, pero no es necesario para la cadena central.
\end{ampliacion}

\section{Aplicaciones y proyección educativa}
Si $N=\sum_{k=0}^{s}a_k10^k$, entonces
\[
N\equiv\sum_{k=0}^{s}a_k r_k\pmod m,\qquad r_k\equiv10^k\pmod m.
\]
Para $m=3$ o $9$, $10^k\equiv1$; para $m=11$, $10^k\equiv(-1)^k$. Las últimas cifras se controlan módulo $10^s$ y los exponentes grandes mediante ciclos o Euler. Bézout resuelve ecuaciones diofánticas y el TCR organiza condiciones simultáneas.

\begin{ampliacion}[· RSA]
En RSA se eligen primos $p,q$, se toma $n=pq$ y exponentes $e,d$ con $ed\equiv1\pmod{\varphi(n)}$. Euler--Fermat, junto con el TCR, justifica la recuperación del mensaje.
\end{ampliacion}

Una secuencia de aula de alto rendimiento consiste en explorar sumas alternadas de cifras, formular el criterio de divisibilidad por $11$ y demostrarlo a partir de $10\equiv-1\pmod{11}$. Se pasa de regularidades a conjetura y prueba. Goñi y colaboradores subrayan que el sentido numérico implica elegir, estimar y justificar estrategias; la tecnología puede explorar casos, pero no sustituye la argumentación.

Las órdenes andaluzas de 30 de mayo de 2023 para ESO y Bachillerato sitúan estos contenidos en el sentido numérico y los conectan con resolución de problemas, razonamiento y pensamiento computacional.

\section{Conclusión}
La división euclídea aporta el principio algorítmico; Bézout caracteriza el MCD mediante combinaciones lineales y resuelve ecuaciones diofánticas; la factorización prima describe la estructura multiplicativa; y las congruencias conservan la información residual compatible con las operaciones. Esta cadena unifica criterios de divisibilidad, sistemas de restos, ciclos de potencias y aplicaciones, y constituye un terreno especialmente fértil para pasar de la experimentación a la prueba.

\clearpage
\section*{Anexo editorial · referencias y control de módulos}
\begin{multicols}{2}\small
\textbf{Bibliografía utilizada}
\begin{itemize}
\item Boyer, C. B. y Merzbach, U. C. (1999). \emph{Historia de la matemática}. Alianza.
\item Hunacek, M. (2023). \emph{Introduction to Number Theory}. CRC.
\item Ivorra Castillo, C. \emph{Introducción a la teoría algebraica de números}. UV.
\item Goñi, J. M. et al. (2011). \emph{Matemáticas. Complementos de formación disciplinar}. Graó.
\item Apostol, T. M. (1984). \emph{Introducción a la teoría analítica de números}. Solo si se conserva la ampliación de funciones aritméticas.
\end{itemize}
\columnbreak
\textbf{Normativa}
\begin{itemize}
\item Orden de 30 de mayo de 2023, currículo de ESO en Andalucía.
\item Orden de 30 de mayo de 2023, currículo de Bachillerato en Andalucía.
\end{itemize}
\end{multicols}

\begin{control}[Tabla de decisión antes del ensayo]
\begin{tabularx}{\textwidth}{L{4.1cm}C{2.3cm}Y}
\toprule
\textbf{Módulo C} & \textbf{Ahorro} & \textbf{Regla de supresión}\\\midrule
Construcción formal de $\Z$ & 4--5 min & Comenzar con la estructura algebraica y división euclídea.\\
Infinitud de los primos & 3--4 min & Conservar lema de Euclides y TFA.\\
Funciones $\tau$ y $\sigma$ & 2--3 min & Retirar también la cita de Apostol.\\
Wilson & 2 min & Suprimir sin referencias posteriores.\\
RSA & 3--4 min & Conservar divisibilidad por 11 como aplicación.\\
\bottomrule
\end{tabularx}
\end{control}

\begin{idea}[Comprobación de cierre]
Antes de considerar esta versión canónica debe escribirse a mano. Registrar: tiempo de cada bloque, páginas, hipótesis omitidas, conclusión terminada y margen real de revisión. La paginación tipográfica no predice la extensión manuscrita.
\end{idea}
\end{document}
